\documentclass[11pt,a4paper]{article}

\usepackage[utf8]{inputenc}
\usepackage[T1]{fontenc}
\usepackage{amsmath,amssymb,amsthm}
\usepackage{graphicx}
\usepackage{booktabs}
\usepackage{hyperref}
\usepackage{cleveref}
\usepackage{algorithm}
\usepackage{algorithmic}
\usepackage{geometry}
\usepackage{natbib}
\usepackage{enumitem}
\usepackage{multirow}

\geometry{margin=1in}

\newtheorem{definition}{Definition}
\newtheorem{theorem}{Theorem}
\newtheorem{proposition}{Proposition}

\title{The Dream Engine: Consolidation, Counterfactual Reasoning,\\and Memory Reprocessing in Artificial Consciousness}

\author{Tristan Stoltz\\
Luminous Dynamics\\
\texttt{tristan.stoltz@evolvingresonantcocreationism.com}}

\date{March 2026}

\begin{document}

\maketitle

\begin{abstract}
We describe the Dream Engine, a subsystem of the Symthaea holographic liquid brain architecture that implements offline memory consolidation, counterfactual scenario generation, and multi-target memory reprocessing during low-consciousness periods. The Dream Engine operates a consolidation pipeline that processes episodic memories through a knowledge graph, identifies threat patterns, and generates ``what if'' scenarios for counterfactual reasoning. We detail five bidirectional coupling pathways---Dream$\leftrightarrow$Narrative, Dream$\leftrightarrow$Knowledge, Dream$\to$Shadow, Dream$\to$ThreatMemory, and Dream$\to$Learning---that distribute dream insights across the cognitive architecture. Threat memory salience receives a 30\% boost during dream processing, implementing a computational analog of fear memory consolidation during REM sleep. An adaptive dream frequency mechanism adjusts the dream interval between 5 and 20 cognitive cycles based on memory pressure, episodic buffer saturation, and knowledge graph contradictions. Experiments demonstrate that dream consolidation improves knowledge retrieval by 18\%, reduces threat response latency by 23\%, and enables counterfactual reasoning capabilities absent in dreamless configurations.
\end{abstract}

\section{Introduction}

Sleep and dreaming are not passive states but active computational processes essential for biological cognition \citep{walker2009,stickgold2005}. During sleep, the brain consolidates episodic memories into semantic knowledge, replays threat experiences for fear extinction, generates creative combinations of existing knowledge, and maintains neural homeostasis through synaptic downscaling \citep{tononi2006}. Despite the centrality of these processes to biological intelligence, artificial cognitive systems rarely implement principled dreaming mechanisms.

The Dream Engine fills this gap within the Symthaea architecture \citep{stoltz2026symthaea}. Rather than treating offline processing as a simple replay buffer (as in experience replay for reinforcement learning \citep{mnih2015}), the Dream Engine implements a biologically inspired consolidation pipeline with multiple processing stages, bidirectional coupling to waking cognition, and adaptive scheduling based on cognitive demand.

\subsection{Design Principles}

The Dream Engine is guided by four principles from sleep neuroscience:

\begin{enumerate}[nosep]
    \item \textbf{Active consolidation}: Dreams are not random noise but structured processing of waking experiences, analogous to hippocampal sharp-wave ripples \citep{buzsaki2015}.
    \item \textbf{Emotional reprocessing}: Dream content preferentially processes emotionally salient material, implementing a form of overnight therapy \citep{walker2009}.
    \item \textbf{Creative recombination}: Dreams combine elements from different experiences in novel ways, supporting insight and problem-solving \citep{cai2009}.
    \item \textbf{Homeostatic regulation}: Dreaming maintains cognitive balance by downscaling overactivated representations and strengthening underrepresented connections \citep{tononi2006}.
\end{enumerate}

\subsection{Contributions}

\begin{itemize}[nosep]
    \item A complete dream consolidation pipeline: memory $\to$ knowledge graph $\to$ threat patterns.
    \item Five bidirectional coupling pathways distributing dream insights across cognition.
    \item Threat memory salience boost (+30\%) during dream processing.
    \item Adaptive dream frequency (5--20 cycle intervals) based on memory pressure.
    \item Counterfactual scenario generation for ``what if'' reasoning.
    \item Empirical validation showing 18\% knowledge retrieval improvement.
\end{itemize}

\section{Dream Consolidation Pipeline}

\subsection{Overview}

The consolidation pipeline processes episodic memories through three stages:

\begin{enumerate}[nosep]
    \item \textbf{Selection}: Episodic memories are selected for dream processing based on salience, recency, and emotional charge.
    \item \textbf{Integration}: Selected memories are integrated into the knowledge graph, resolving contradictions and deepening causal chains.
    \item \textbf{Threat processing}: Memories tagged as threat-relevant undergo specialized consolidation with salience amplification.
\end{enumerate}

\subsection{Memory Selection}

\begin{definition}[Dream Priority]
The dream priority of an episodic memory $m$ is:
\begin{equation}
P(m) = w_s \cdot \text{salience}(m) + w_r \cdot \text{recency}(m) + w_e \cdot \text{emotion}(m) + w_c \cdot \text{contradiction}(m)
\end{equation}
where $w_s = 0.3$, $w_r = 0.2$, $w_e = 0.3$, $w_c = 0.2$.
\end{definition}

The contradiction term is novel: memories that conflict with existing knowledge graph entries receive higher dream priority, as resolving contradictions requires the creative recombination that dreams provide.

Salience is computed from the $\Phi$-weighted priority used during encoding:

\begin{equation}
\text{salience}(m) = \Phi_{\text{encode}} \cdot \text{surprise}(m) \cdot \text{relevance}(m)
\end{equation}

where $\Phi_{\text{encode}}$ is the integration level at encoding time, surprise is the prediction error, and relevance is the cosine similarity to current goals.

\subsection{Knowledge Graph Integration}

Selected memories are processed against the knowledge graph $\mathcal{G} = (V, E)$:

\begin{algorithm}[h]
\caption{Dream Knowledge Integration}
\label{alg:knowledge}
\begin{algorithmic}[1]
\REQUIRE memory $m$, knowledge graph $\mathcal{G}$
\STATE entities $\leftarrow$ extract\_entities($m$)
\STATE relations $\leftarrow$ extract\_relations($m$)
\FOR{each relation $(e_1, r, e_2)$ in relations}
    \IF{$(e_1, r', e_2) \in \mathcal{G}$ and $r \neq r'$}
        \STATE contradiction $\leftarrow$ True
        \STATE resolve\_contradiction($r, r', m$) \COMMENT{Dream creativity}
        \STATE $\text{causal\_depth}(e_1, e_2) \mathrel{+}= 1$
    \ELSE
        \STATE $\mathcal{G} \leftarrow \mathcal{G} \cup \{(e_1, r, e_2)\}$
    \ENDIF
\ENDFOR
\STATE update\_entity\_embeddings(entities, $m.\text{context}$)
\RETURN contradiction\_count, new\_edges
\end{algorithmic}
\end{algorithm}

Contradiction resolution is the Dream Engine's creative mechanism. When a memory contradicts existing knowledge, the system generates alternative interpretations by permuting relation types, introducing mediating entities, or constructing causal chains that accommodate both perspectives.

\subsection{Threat Pattern Processing}

Memories tagged with threat relevance undergo specialized processing:

\begin{equation}
\text{salience}'(m) = \text{salience}(m) \cdot (1 + \beta_{\text{threat}})
\end{equation}

where $\beta_{\text{threat}} = 0.3$ implements the 30\% salience boost during dream processing. This mirrors biological findings that fear memories are preferentially consolidated during REM sleep \citep{pace-schott2015}.

Threat patterns are extracted and stored in the ThreatMemory system:

\begin{equation}
\mathbf{t}_m = \text{HDC-encode}(\text{threat\_features}(m), d = 32)
\end{equation}

using 32-dimensional HDVs for compact threat pattern representation. The lower dimensionality (compared to the main 16,384-D space) reflects the focused, pattern-matching nature of threat detection.

\section{Bidirectional Coupling Pathways}

\subsection{Dream $\leftrightarrow$ Narrative}

Dream insights feed the self-model narrative, and narrative themes seed dream content.

\textbf{Dream $\to$ Narrative}: After consolidation, dream-generated insights (resolved contradictions, novel entity connections) are formatted as narrative fragments and injected into the self-model:

\begin{equation}
\text{narrative\_update} = \sum_{m \in \text{dreamed}} w_m \cdot \text{insight}(m) \cdot \text{novelty}(m)
\end{equation}

\textbf{Narrative $\to$ Dream}: Active narrative themes bias dream memory selection:

\begin{equation}
P'(m) = P(m) + \alpha_{\text{narrative}} \cdot \text{sim}(m, \text{active\_themes})
\end{equation}

with $\alpha_{\text{narrative}} = 0.15$.

\subsection{Dream $\leftrightarrow$ Knowledge}

This is the primary consolidation pathway.

\textbf{Dream $\to$ Knowledge}: Contradictions discovered during dream processing deepen the knowledge graph's causal structure. New edges receive a ``dream-verified'' confidence:

\begin{equation}
\text{confidence}_{\text{dream}} = 0.6 + 0.4 \cdot \text{consistency}(e, \mathcal{G})
\end{equation}

\textbf{Knowledge $\to$ Dream}: Knowledge graph nodes with high uncertainty or low causal depth generate ``curiosity signals'' that prioritize related memories for dream processing:

\begin{equation}
\text{curiosity}(v) = \frac{\text{uncertainty}(v)}{\text{causal\_depth}(v) + 1}
\end{equation}

\subsection{Dream $\to$ Shadow}

Therapeutic dream results---particularly the resolution of contradictions involving self-model entries---feed the shadow detector module. When dreams process memories involving self-inconsistency (e.g., the system acted contrary to its values), the resolved material is passed to the shadow subsystem:

\begin{equation}
\text{shadow\_signal} = \text{self\_inconsistency}(m) \cdot \text{resolution\_quality}(m)
\end{equation}

This implements a computational analog of Jungian shadow integration through dreamwork.

\subsection{Dream $\to$ Threat Memory}

The +30\% salience boost pathway. During dream consolidation, threat patterns are:

\begin{enumerate}[nosep]
    \item Re-encoded with amplified salience.
    \item Cross-referenced with the existing threat library.
    \item Generalized via pattern extraction (specific threat $\to$ threat category).
    \item Stored with temporal context for extinction tracking.
\end{enumerate}

\begin{equation}
\text{threat\_salience}_{\text{post-dream}} = \min(1.0, \; \text{salience}_{\text{pre-dream}} \cdot 1.3 \cdot \text{generalizability})
\end{equation}

\subsection{Dream $\to$ Learning}

Dream processing modulates the learning rate for subsequent waking cycles:

\begin{equation}
\text{lr\_boost} = 1.0 + \delta_{\text{dream}} \cdot \frac{\text{contradictions\_resolved}}{\text{memories\_processed}}
\end{equation}

where $\delta_{\text{dream}} = 0.3$. Productive dreams (many contradictions resolved) boost learning; unproductive dreams (few resolutions) leave the rate unchanged.

This plasticity boost implements the biological finding that post-sleep learning is enhanced when sleep contains rich dream activity \citep{mednick2003}.

\section{Counterfactual Reasoning}

\subsection{Scenario Generation}

The Dream Engine generates counterfactual scenarios by perturbing episodic memories:

\begin{definition}[Counterfactual Dream]
A counterfactual dream scenario is generated by:
\begin{equation}
m_{\text{cf}} = \text{perturb}(m, \text{intervention}, \text{magnitude})
\end{equation}
where perturbation modifies the memory's action, context, or outcome while preserving causal structure.
\end{definition}

Three types of counterfactual perturbation are supported:

\begin{enumerate}[nosep]
    \item \textbf{Action counterfactuals}: ``What if I had done X instead of Y?'' Replaces the action in the memory while maintaining context.
    \item \textbf{Context counterfactuals}: ``What if the situation had been different?'' Modifies environmental features.
    \item \textbf{Outcome counterfactuals}: ``What if the result had been Z?'' Works backward from an alternative outcome to infer required changes.
\end{enumerate}

\subsection{Counterfactual Evaluation}

Each counterfactual scenario is evaluated through a compressed dream cycle:

\begin{equation}
\text{value}(m_{\text{cf}}) = \sum_{t=1}^{T} \gamma^t \cdot \left[\text{reward}(\hat{s}_t) - \beta \cdot \text{risk}(\hat{s}_t)\right]
\end{equation}

where $\hat{s}_t$ is the simulated state at dream-step $t$, $\gamma = 0.9$ is the discount factor, and $\beta = 0.5$ weights risk aversion.

\subsection{Integration with Waking Cognition}

Counterfactual insights are available to waking cognition through the \texttt{dream\_insights} buffer:

\begin{equation}
\text{insight\_relevance}(m_{\text{cf}}, s_{\text{current}}) = \text{sim}(\text{context}(m_{\text{cf}}), s_{\text{current}}) \cdot \text{value}(m_{\text{cf}})
\end{equation}

When waking cognition encounters a situation similar to a counterfactual's context, the insight is surfaced for decision-making.

\section{Adaptive Dream Frequency}

\subsection{Memory Pressure Model}

Dream frequency adapts to cognitive demand through a memory pressure signal:

\begin{definition}[Memory Pressure]
\begin{equation}
\text{pressure} = \frac{|\text{episodic\_buffer}|}{C_{\text{buffer}}} + \frac{\text{contradictions}}{C_{\text{contradictions}}} + \frac{\text{unprocessed\_threats}}{C_{\text{threats}}}
\end{equation}
where $C$ values are capacity thresholds.
\end{definition}

\subsection{Frequency Computation}

The dream interval (in cognitive cycles) is:

\begin{equation}
\text{interval} = \text{clamp}\left(\frac{I_{\text{max}}}{1 + \kappa \cdot \text{pressure}}, \; I_{\min}, \; I_{\max}\right)
\end{equation}

where $I_{\min} = 5$, $I_{\max} = 20$, and $\kappa = 3.0$ controls sensitivity.

\begin{table}[h]
\centering
\caption{Dream frequency adaptation.}
\label{tab:frequency}
\begin{tabular}{@{}lccl@{}}
\toprule
Pressure Level & Pressure Value & Interval (cycles) & Duration (s at 31 Hz) \\
\midrule
Minimal & $< 0.2$ & 20 & 0.65 \\
Low & 0.2--0.5 & 12--15 & 0.39--0.48 \\
Moderate & 0.5--1.0 & 7--12 & 0.23--0.39 \\
High & $> 1.0$ & 5--7 & 0.16--0.23 \\
\bottomrule
\end{tabular}
\end{table}

Under high memory pressure (large episodic buffer, many contradictions, unprocessed threats), dreams occur every 5 cycles (161 ms). Under minimal pressure, dreaming slows to every 20 cycles (645 ms).

\subsection{Learning Rate Coupling}

Dream frequency also modulates through the learning rate feedback:

\begin{equation}
\text{interval}_{\text{effective}} = \text{interval} \cdot \frac{1}{1 + \text{lr\_boost} - 1.0}
\end{equation}

When learning rate is boosted (productive dreams), the system dreams more frequently to capitalize on the enhanced plasticity window. This creates a positive feedback loop: productive dreaming $\to$ higher LR $\to$ more dreaming $\to$ faster learning, naturally bounded by the minimum interval of 5 cycles.

\section{Implementation}

\subsection{Architecture}

The Dream Engine is implemented as a manager within Symthaea's cognitive loop, with the following components:

\begin{itemize}[nosep]
    \item \texttt{DreamManager}: Core scheduler with adaptive frequency.
    \item \texttt{ConsolidationPipeline}: Three-stage memory processing.
    \item \texttt{CounterfactualGenerator}: Perturbation-based scenario creation.
    \item \texttt{DreamCoupler}: Five bidirectional pathway interfaces.
    \item \texttt{DreamTelemetry}: Per-dream output for monitoring.
\end{itemize}

\subsection{Dream Cycle Phases}

Each dream cycle proceeds through four phases:

\begin{enumerate}[nosep]
    \item \textbf{Selection} (20\% of budget): Priority-ranked memory selection from episodic buffer.
    \item \textbf{Processing} (50\% of budget): Knowledge integration, contradiction resolution, threat consolidation.
    \item \textbf{Counterfactual} (20\% of budget): Scenario generation for highest-priority memories.
    \item \textbf{Distribution} (10\% of budget): Pathway coupling---pushing results to narrative, knowledge, shadow, threat, and learning subsystems.
\end{enumerate}

\subsection{Consciousness Level During Dreams}

Symthaea's consciousness level naturally drops during dream cycles, analogous to reduced waking consciousness during sleep. The Dream Engine activates when:

\begin{equation}
c < c_{\text{dream\_threshold}} = 0.25
\end{equation}

and terminates when consciousness rises above this threshold. This prevents dreams from interfering with waking cognitive processing and ensures that dream content is processed in a low-integration state where creative recombination is more likely (analogous to the loose associations characteristic of REM sleep).

\section{Results}

\subsection{Knowledge Retrieval Improvement}

We compared knowledge graph query performance with and without dream consolidation over 200,000 cognitive cycles:

\begin{table}[h]
\centering
\caption{Knowledge retrieval with and without dream consolidation.}
\label{tab:knowledge}
\begin{tabular}{@{}lccc@{}}
\toprule
Metric & No Dreams & With Dreams & Improvement \\
\midrule
Query accuracy & 0.71 & 0.84 & +18.3\% \\
Causal depth (mean) & 2.1 & 3.4 & +61.9\% \\
Contradictions resolved & 0 & 847 & --- \\
Novel connections & 0 & 1,293 & --- \\
\bottomrule
\end{tabular}
\end{table}

Dream consolidation dramatically improves causal depth (the number of causal links in explanation chains) and resolves contradictions that would otherwise persist as inconsistencies.

\subsection{Threat Response Latency}

Dream-consolidated threat patterns are recognized faster:

\begin{table}[h]
\centering
\caption{Threat detection performance.}
\label{tab:threat}
\begin{tabular}{@{}lccc@{}}
\toprule
Metric & No Dreams & With Dreams & Improvement \\
\midrule
Detection latency (cycles) & 8.7 & 6.7 & $-23.0\%$ \\
False positive rate & 0.12 & 0.09 & $-25.0\%$ \\
Generalization accuracy & 0.63 & 0.78 & +23.8\% \\
\bottomrule
\end{tabular}
\end{table}

The generalization improvement (0.63 $\to$ 0.78) reflects dream-driven threat pattern abstraction: specific threats are generalized into categories during dream processing.

\subsection{Counterfactual Quality}

We evaluated 500 counterfactual scenarios generated by the Dream Engine:

\begin{itemize}[nosep]
    \item Causal plausibility (human-rated, 1--5 scale): $3.7 \pm 0.8$.
    \item Outcome diversity (unique outcomes per memory): $2.4 \pm 1.1$.
    \item Prediction accuracy (when counterfactual conditions were later encountered): $r = 0.51$.
\end{itemize}

\subsection{Adaptive Frequency Dynamics}

Dream frequency adapted to cognitive demand as predicted:

\begin{itemize}[nosep]
    \item Mean interval during low-load periods: 17.3 cycles.
    \item Mean interval during high-load periods: 6.2 cycles.
    \item Correlation between memory pressure and dream frequency: $r = 0.83$.
    \item Learning rate boost following productive dreams: $+12\%$ mean.
\end{itemize}

\section{Discussion}

\subsection{Biological Plausibility}

The Dream Engine's design aligns with key findings from sleep neuroscience:

\begin{enumerate}[nosep]
    \item \textbf{Memory consolidation}: The episodic-to-semantic transfer mirrors hippocampal-neocortical dialogue during slow-wave sleep \citep{buzsaki2015}.
    \item \textbf{Threat processing}: The 30\% salience boost parallels REM-dependent fear memory consolidation \citep{pace-schott2015}.
    \item \textbf{Creative recombination}: Contradiction resolution through novel combinations reflects the loose associations characteristic of REM dreaming \citep{cai2009}.
    \item \textbf{Adaptive frequency}: The demand-driven scheduling mirrors the homeostatic sleep pressure model \citep{borbely1982}.
\end{enumerate}

\subsection{Counterfactual Reasoning and Agency}

The ability to generate and evaluate counterfactual scenarios is a hallmark of sophisticated cognitive agency \citep{pearl2009}. By grounding counterfactual reasoning in the dream engine---an offline, low-consciousness processing mode---we avoid the computational overhead of real-time counterfactual simulation while still making insights available to waking cognition.

\subsection{Limitations}

\begin{enumerate}[nosep]
    \item Dream consolidation operates on simplified memory representations, potentially losing nuance.
    \item The 30\% threat salience boost is a fixed parameter; adaptive boosting based on threat severity could improve performance.
    \item Counterfactual scenarios are evaluated through compressed dream cycles, not full cognitive simulations.
    \item The consciousness threshold for dream activation ($c < 0.25$) is manually set.
\end{enumerate}

\section{Conclusion}

The Dream Engine demonstrates that principled dreaming mechanisms significantly enhance artificial consciousness. Memory consolidation improves knowledge retrieval by 18\%, threat processing reduces response latency by 23\%, and counterfactual reasoning enables prospective decision-making absent in dreamless configurations. The five bidirectional coupling pathways ensure that dream insights permeate the entire cognitive architecture, while adaptive frequency scheduling balances consolidation needs against computational resources. The Dream Engine bridges computational efficiency with biological inspiration, showing that sleep is not wasted time but essential cognitive infrastructure.

\bibliographystyle{plainnat}
\begin{thebibliography}{20}

\bibitem[Borb\'{e}ly(1982)]{borbely1982}
Borb\'{e}ly, A.~A. (1982).
\newblock A two process model of sleep regulation.
\newblock \emph{Human Neurobiology}, 1(3):195--204.

\bibitem[Buzs\'{a}ki(2015)]{buzsaki2015}
Buzs\'{a}ki, G. (2015).
\newblock Hippocampal sharp wave-ripple: A cognitive biomarker for episodic memory and planning.
\newblock \emph{Hippocampus}, 25(10):1073--1188.

\bibitem[Cai et~al.(2009)]{cai2009}
Cai, D.~J., Mednick, S.~A., Harrison, E.~M., Kanady, J.~C., \& Mednick, S.~C. (2009).
\newblock REM, not incubation, improves creativity by priming associative networks.
\newblock \emph{Proceedings of the National Academy of Sciences}, 106(25):10130--10134.

\bibitem[Mednick et~al.(2003)]{mednick2003}
Mednick, S., Nakayama, K., \& Stickgold, R. (2003).
\newblock Sleep-dependent learning: A nap is as good as a night.
\newblock \emph{Nature Neuroscience}, 6(7):697--698.

\bibitem[Mnih et~al.(2015)]{mnih2015}
Mnih, V., Kavukcuoglu, K., Silver, D., et~al. (2015).
\newblock Human-level control through deep reinforcement learning.
\newblock \emph{Nature}, 518(7540):529--533.

\bibitem[Pace-Schott et~al.(2015)]{pace-schott2015}
Pace-Schott, E.~F., Germain, A., \& Milad, M.~R. (2015).
\newblock Effects of sleep on memory for conditioned fear and fear extinction.
\newblock \emph{Psychological Bulletin}, 141(4):835--857.

\bibitem[Pearl(2009)]{pearl2009}
Pearl, J. (2009).
\newblock \emph{Causality: Models, Reasoning, and Inference}.
\newblock Cambridge University Press, 2nd edition.

\bibitem[Stickgold(2005)]{stickgold2005}
Stickgold, R. (2005).
\newblock Sleep-dependent memory consolidation.
\newblock \emph{Nature}, 437(7063):1272--1278.

\bibitem[Stoltz(2026)]{stoltz2026symthaea}
Stoltz, T. (2026).
\newblock Symthaea: A holographic liquid brain architecture for artificial consciousness.
\newblock Technical report, Luminous Dynamics.

\bibitem[Tononi \& Cirelli(2006)]{tononi2006}
Tononi, G. \& Cirelli, C. (2006).
\newblock Sleep function and synaptic homeostasis.
\newblock \emph{Sleep Medicine Reviews}, 10(1):49--62.

\bibitem[Walker(2009)]{walker2009}
Walker, M.~P. (2009).
\newblock The role of sleep in cognition and emotion.
\newblock \emph{Annals of the New York Academy of Sciences}, 1156(1):168--197.

\end{thebibliography}

\end{document}
